\section{Experiments and Results}
\label{sec:experiments}

\subsection{Vision quality across conditions}
Table~\ref{tab:miou} reports the current 100k-step seed-0 run. mIoU is monotonic in model size at every condition. ACDC/night is the hardest condition. The fog score being higher than clean Cityscapes for some levels is a real, investigated result: structural classes improve in visually simpler fog scenes while dynamic road-user classes degrade, so per-class results are required alongside macro mIoU.

\begin{table}[t]
\caption{Current validation mIoU at 100k steps (seed 0).}
\label{tab:miou}
\centering
\begin{tabular}{lrrrrr}
\toprule
Level & Cityscapes & Fog & Night & Rain & Snow \\
\midrule
\texttt{tiny} & 0.2986 & 0.3154 & 0.1944 & 0.2901 & 0.2610 \\
\texttt{small} & 0.3564 & 0.3658 & 0.2511 & 0.3680 & 0.3304 \\
\texttt{medium} & 0.4120 & 0.4360 & 0.2898 & 0.4003 & 0.3981 \\
\texttt{large} & 0.4712 & 0.4919 & 0.3312 & 0.4517 & 0.4492 \\
\bottomrule
\end{tabular}
\end{table}

\subsection{Shared supernet versus an independent baseline}
The corrected apples-to-apples augmented comparison against Fast-SCNN is shown in Table~\ref{tab:baseline}. The supernet wins four of five splits by 0.3--1.3 mIoU points and loses ACDC/rain by 1.2 points. This first pair uses one seed per side; the margins are smaller than observed seed-to-seed variation and must not be presented as a definitive multi-seed result.

\begin{table}[t]
\caption{First augmented comparison with Fast-SCNN (one seed per side).}
\label{tab:baseline}
\centering
\begin{tabular}{lrrr}
\toprule
Dataset & Fast-SCNN & Supernet-large & Gap \\
\midrule
Cityscapes & 0.5227 & 0.5354 & +0.0127 \\
ACDC/fog & 0.5730 & 0.5761 & +0.0031 \\
ACDC/night & 0.3813 & 0.3870 & +0.0057 \\
ACDC/rain & 0.5004 & 0.4883 & -0.0121 \\
ACDC/snow & 0.5055 & 0.5163 & +0.0108 \\
\bottomrule
\end{tabular}
\end{table}

\subsection{Hardware-dependent Pareto choice}
Joining measured latency with the quality table makes all four levels Pareto-optimal on the two currently benchmarked devices. Under a common 10 ms budget, the lookup search selects \texttt{small} on Hailo-8 and \texttt{large} on the TensorRT GPU. This is concrete evidence that a one-size-fits-all subnet choice is not interchangeable across devices, but it is not yet the planned quantitative comparison against a FLOPs-aware baseline.

\TODO{Add three-seed confidence intervals, the remaining baselines, FLOPs-aware selection, E2/E5 data, and complete statistical tests before citing headline claims.}
